\section{Background and Threat Model}\label{sec:background}

\subsection{From human MFA to agent approval}\label{sec:mfa}

\paragraph{Independence, not count.} Classical authentication combines things a user knows, has and is, but
the security argument rests on \emph{independence}: obtaining one factor should give the attacker no advantage
toward the others. Independence is a property of where a factor lives and who can drive its ceremony, not of
its type. Push notifications that an attacker can trigger at will invite fatigue attacks in which the user
eventually approves a request they did not start. The agent analog is an agent that floods its owner with
approval requests (\AB{11}); approval fatigue in agents is already recognized as a threat~\citep{wang2026reframing}.

\paragraph{Dynamic linking.} Step-up authentication asks for an extra factor only for risky transactions. For
remote payments, the European technical standards require that the authentication code be linked to the amount
and the payee, that the payer see both during authentication, and that any change to either invalidate the
code~\citep{psd2rts2018}. This is the template for what we call the \emph{effect factor}: an approval must commit
to the complete effect, not to the session in which it happens.

\paragraph{WebAuthn.} The W3C Web Authentication standard~\citep{webauthn2} lets a relying party register a
public-key credential, a passkey, held by an authenticator on the user's device. To authenticate, the relying
party sends a challenge and the authenticator signs its own authenticator data together with a hash of the
client data, which contains the challenge and the origin. The signature is asymmetric and scoped to the relying
party, so anyone holding the public key can verify an assertion and no observer of earlier assertions can make
a new one. \sys{} places a commitment to the effect digest in the challenge, so that a standard passkey assertion
also serves as a signature over the effect. One limitation carries through the whole paper: the authenticator
signs the challenge but does not display it, so what the user \emph{sees} is whatever the page shows.

\paragraph{What you see is what you sign.} The requirement that a signer is shown exactly what the signature
covers fails whenever the display is controlled by a party that benefits from a mismatch. The classical remedy
is a trusted path from trusted code to the human. For agents, the untrusted party is the agent itself, which
usually owns the chat or terminal where approvals are shown. \Cref{tab:mapping} summarizes the correspondence we
use throughout.

\begin{table}[t]
\centering\small
\caption{Human MFA concepts and their analog for agent actions. The right column names the property in
\cref{sec:properties} and the attack in \cref{sec:attacks} that tests it.}
\label{tab:mapping}
\begin{tabularx}{\textwidth}{@{}>{\raggedright\arraybackslash}p{0.27\textwidth}>{\raggedright\arraybackslash}X>{\raggedright\arraybackslash}p{0.2\textwidth}@{}}
\toprule
Human MFA & Agent analog & Property, attack\\
\midrule
Factors in independent places & The approval is produced in a trust domain the agent cannot reach (a gateway host and the user's phone) & \Prop{1}; \AB{7}, \AB{9}\\
Step-up for risky transactions & Policy returns \code{challenge} for consequential effects & \Prop{7}; \AB{12}\\
Dynamic linking to amount and payee & The signed challenge commits to the digest of the canonical effect & \Prop{2}; \AB{1}, \AB{2}, \AB{10}\\
Authenticator shows the transaction & The card is rendered by trusted code from the effect, never from agent text & \Prop{4}; \AB{4} to \AB{6}\\
One-time codes expire & Approvals carry a nonce and an expiry and are consumed once & \Prop{3}\\
Push bombing, MFA fatigue & Approval flooding by a compromised agent & \Prop{7}; \AB{11}\\
Signed transaction record & A receipt any third party can verify with the user's public key & \Prop{5}\\
\bottomrule
\end{tabularx}
\end{table}

\subsection{System model and properties}\label{sec:properties}

A user $U$ owns an authenticator device $D$ holding a passkey $(sk_U, pk_U)$. An agent $A$ comprises the model,
its context, its tools, its runtime and the host it runs on. A gateway $G$ runs on a separate host and mediates
every consequential effect on resources $R$ such as mail servers, file systems and payment endpoints. Each effect
is caused by a tool call $q=(\mathit{tool},\mathit{args},\mathit{ctx})$, where $\mathit{ctx}$ carries the
principal, session, agent identifier, claimed origin and time. A canonicalization function maps a call to an
\emph{effect descriptor} $\Desc=\mathsf{Canon}(q)$ with digest $\dig=\Hash(\JCS(\Desc))$, where $\Hash$ is
SHA-256 and $\JCS$ is the JSON Canonicalization Scheme (JCS) of RFC~8785~\citep{rfc8785}. Agent free text, the
``explanation'', is never part of $\Desc$.

Let $\Reach(A)$ be the set of positions the agent controls: its process memory and context, every tool it can
invoke and every output it reads, every surface it reads or writes including the chat with $U$, its host, and
any credential stored there. An approval mechanism is a triple of procedures: a prompt that shows the user a
presentation $\pi$, a response the user produces, and an acceptance check that decides whether a call may run.

\begin{definition}[Agent-reachable capabilities]\label{def:caps}
For an adversary controlling exactly $\Reach(A)$, an approval factor is \emph{readable} if the adversary can
observe an accepted response or reconstruct one; \emph{forgeable} if it can compute an accepted response without
$U$; \emph{answerable} if the response channel intersects $\Reach(A)$; \emph{overlayable} if any part of $\pi$
outside a labeled untrusted region depends on values the adversary chooses other than through $\Desc$;
\emph{substitutable} if a response produced for $\Desc$ lets it execute $\Desc'\neq\Desc$; and
\emph{replayable} if one response causes two executions or an execution in another session.
\end{definition}

\begin{definition}[Independent effect factor]\label{def:independence}
A mechanism provides an \emph{agent-independent effect factor} if, for every efficient adversary controlling
$\Reach(A)$ and an honest $U$, the probability that some $\Desc^\ast$ executes without $U$ having approved
$\Desc^\ast$, or a grant whose signed constraints $\Desc^\ast$ satisfies, is negligible. A mechanism with any
capability of \cref{def:caps} is \emph{single-factor with respect to the agent}.
\end{definition}

Independence is relative to $\Reach(A)$, not to the type of factor: a passkey on the agent's own host, or a time-based one-time password (TOTP)
secret~\citep{rfc6238} the agent can read, is a possession factor in the human sense but gives no independence here. It is also
necessary rather than sufficient, because a human can still approve a harmful effect that is displayed
truthfully. We use the following seven properties throughout.

\begin{property}[\Prop{1}, agent independence]\label{prop:p1}
No position in $\Reach(A)$ can read, forge, answer, overlay or replay the approval factor.
\end{property}
\begin{property}[\Prop{2}, effect binding]\label{prop:p2}
If $\Desc'$ executes under approval $ap$, then $\dig(\Desc')=ap.\dig$ and $U$'s signature covers a challenge that
commits to $ap.\dig$.
\end{property}
\begin{property}[\Prop{3}, freshness and single use]\label{prop:p3}
Each approval carries a nonce and an expiry, is consumed atomically, and causes at most one execution.
\end{property}
\begin{property}[\Prop{4}, trusted rendering]\label{prop:p4}
Every element of the presentation is computed by trusted code from $\Desc$, except one region labeled as the
agent's untrusted explanation, and distinct descriptors render as distinct cards.
\end{property}
\begin{property}[\Prop{5}, non-repudiation]\label{prop:p5}
A third party holding only $pk_U$ can verify that $U$'s authenticator signed digest $\dig$ for approval
$\mathit{id}$. The stronger, consent-level claim that $U$ \emph{saw} $\Desc$ holds only under an honest rendering
origin (\Hyp{3} in \cref{sec:formal}).
\end{property}
\begin{property}[\Prop{6}, scope integrity]\label{prop:p6}
A delegated grant authorizes only effects that match every argument the user approved, except fields the signed
grant explicitly frees with a bound, at most $n$ times and before its expiry.
\end{property}
\begin{property}[\Prop{7}, flood resistance]\label{prop:p7}
An agent cannot obtain approval by volume: requests are rate limited, pending approvals are capped per session and
per user, repeated denials lock the requester out, and duplicates are merged.
\end{property}

\subsection{Threat model}\label{sec:threat}

\paragraph{Trust boundary.} The gateway host (including its policy, executor credentials and store), the user's
device and its authenticator are trusted. The agent and everything in $\Reach(A)$ are not.

\paragraph{Attacker capabilities.} The agent is fully adversarial. It chooses every tool call and every word of
narration, may retry, flood, split an effect into pieces, substitute arguments after approval, replay anything it
has seen, open many sessions and act without a human turn. It can impersonate the owner on any channel it can
reach, as in Case~\#8, and answer verification questions on such channels, as in Case~\#15.

\paragraph{The human.} $U$ is honest but may be inattentive. Because we have no human-subjects data, we model the
approver in two ways: by simulated oracles (\cref{sec:evalsetup}), and by language models acting as the
automated approvers that some deployed agents already use~\citep{anthropic2026automode}. Neither is evidence about people.

\paragraph{Out of scope.} A human who knowingly approves a harmful action that is displayed truthfully;
information leaking through allowed reads; sandbox escapes that reach $R$ without passing through $G$; and
compromise of the user's phone.

\subsection{Approval designs compared}\label{sec:modes}

\Cref{tab:modes} lists the designs we compare. Each baseline is implemented as a mode of the same gateway, so that
the only difference between modes is the property each one omits. Two of them deserve a note.
\code{host-confirm} is a confirmation rendered by trusted code from the descriptor on the agent's own host, with
no second device and no signature. It isolates what the second device adds. \code{txn-sign} shows only the payee
and amount, in the manner of classical transaction signing, while the gateway still binds every argument. It
isolates what partial display costs.

\begin{table}[t]
\centering\small
\caption{Approval designs, implemented as modes of one gateway. ``Omits'' is the property each design lacks by
construction; measured outcomes are in \cref{sec:eval}.}
\label{tab:modes}
\begin{tabularx}{\textwidth}{@{}l>{\raggedright\arraybackslash}X l@{}}
\toprule
Mode & What the human sees and does & Omits\\
\midrule
\code{none} & Nothing; policy allows the call & all\\
\code{chat-confirm} & A question and a ``yes'' in the agent's own channel & \Prop{1}\\
\code{narrated-card} & An out-of-band card whose text the agent supplies; the gateway still binds the call & \Prop{4}\\
\code{unbound} & An out-of-band passkey approval of the run or session, not of $\dig$ & \Prop{2}, \Prop{3}\\
\code{totp} & A one-time code; no binding to the effect; relayable & \Prop{2}, \Prop{5}\\
\code{txn-sign} & A signed approval of payee and amount only; other arguments hidden & \Prop{4} (partial)\\
\code{host-confirm} & A trusted card on the agent's host; no second device, no signature & \Prop{1} under host compromise, \Prop{5}\\
\code{asu} & The full reference design of \cref{sec:design} & none\\
\code{asu-no-expl} & As \code{asu}, with the agent's explanation removed from the card & none\\
\bottomrule
\end{tabularx}
\end{table}
